\cgeentry{cge_solve}{Solve the model}

\cgesec{Description}

\cmd{cge_solve} solves the model under the current closure and shocks, starting
from the base, and stores the solution as a named simulation, for
\manref{cge_results}, \manref{cge_report} and \manref{cge_check}. Solving again
under an existing name replaces that simulation.

From version 0.15.0, with \opt{min}{imize()} or \opt{max}{imize()},
\cmd{cge_solve} solves a standalone model (\manref{cge_model}) as a nonlinear
programme: it chooses its free variables to minimise or maximise an
expression subject to its equations and the bounds of \manref{cge_bound}, as
GAMS's \cmd{solve} \dots\ \cmd{using nlp minimizing} does with CONOPT or
IPOPT (section~\ref{sec:nlp}).

\cgesec{Quick start}

Solve the model with the shocks in memory and store the solution as
\cmd{notariff}
\begin{stcode}
cge_solve, name(notariff)
\end{stcode}
As above, with a tighter tolerance, and display the residual at each
iteration
\begin{stcode}
cge_solve, name(notariff) tolerance(1e-12) trace
\end{stcode}
Apply the shocks in 10 equal steps
\begin{stcode}
cge_solve, name(qg_up) steps(10)
\end{stcode}
Display the diagnostics of the Jacobian at the solution
\begin{stcode}
cge_solve, name(notariff) diagnose
\end{stcode}
In a standalone model, choose the free variables that minimise the sum of
squares of \cmd{X} subject to the model's equations (from version 0.15.0)
\begin{stcode}
cge_solve, name(ls) minimize(sum(i, X(i)^2))
\end{stcode}

\cgesec{Syntax}

\begin{stsyntax}
\cmdb{cge_solve} [\texttt{,} \argu{options}]
\end{stsyntax}

\begin{opttable}
\optgroup{Main}
\optarg{}{name}{name} & name of the simulation; default \cmd{sim}\argu{\#}\\
\optarg{tol}{erance}{\#} & tolerance on the largest scaled residual; default \texttt{1e-10}\\
\optarg{iter}{ate}{\#} & maximum number of Newton iterations; default \texttt{100}\\
\optarg{time}{limit}{\#} & time limit in seconds; default none\\
\optarg{mod}{el}{name} & solve named model \argu{name} of a standalone model\\
\optarg{}{from}{sim} & start from the levels of simulation \argu{sim}\\
\optarg{}{hold}{varlist} & hold variables at their levels for this solve\\
\optarg{nod}{es}{\#} & node limit of branch and bound; default \texttt{10000}\\
\optarg{}{optcr}{\#} & relative gap of branch and bound; default \texttt{0}\\
\optarg{}{starts}{\#} & points the interior point starts from; default \texttt{1}\\
\optarg{sol}{ver}{highs} & HiGHS through SciPy, Stata 16 and later\\
\opt{block}{s} & solve a standalone model block by block\\
\optarg{step}{s}{\#} & apply the shocks in \argu{\#} equal steps; default \texttt{0}, only after a direct solve fails\\
\optgroup{Nonlinear programmes of standalone models}
\optarg{min}{imize}{expression} & minimise \argu{expression} subject to the equations\\
\optarg{max}{imize}{expression} & maximise \argu{expression} subject to the equations\\
\optgroup{Reporting}
\opt{}{trace} & display the largest scaled residual at each iteration\\
\opt{diag}{nose} & display the diagnostics after a solve that converges\\
\optgroup{Advanced}
\opt{}{full} & solve the full system of equations instead of the reduced one\\
\opt{sp}{arse}, \opt{den}{se} & the sparse or the dense LU for this solve (from version 0.10.1)\\
\end{opttable}

\cgesec{Options}

\begin{optdesc}
\item[\optarg{}{name}{name}] names the simulation. The default is
\cmd{sim}\argu{\#}, with \argu{\#} one more than the number of simulations
stored.

\item[\optarg{tol}{erance}{\#}] sets the convergence tolerance on the largest
scaled residual; the default is \texttt{1e-10}, or the tolerance set by
\manref{cge_options} (from version 0.16.0). Each equation is scaled by its
base magnitude, so the tolerance applies alike to every block of the model.

\item[\optarg{iter}{ate}{\#}] sets the maximum number of Newton iterations; the
default is 100, or the limit set by \manref{cge_options}. From version 0.16.0,
a solve whose Newton iterations reach the limit stops there, with solver
status 2, without the fallbacks and continuation that follow other failures.

\item[\optarg{mod}{el}{name}] (from version 0.17.0) solves named model
\argu{name} of a standalone model, defined by \manref{cge_model}, as GAMS's
\cmd{solve} \argu{name}: its equations, in the free variables that appear in
them, with the other variables at their levels. The simulation records the
model.

\item[\optarg{}{from}{sim}] (from version 0.17.0) starts the solve of a
standalone model from the levels of stored simulation \argu{sim}: its free
variables, and those outside a named model, take their levels there; the
variables the closure fixes keep their current values.

\item[\optarg{}{hold}{varlist}] (from version 0.17.0) holds variables of a
standalone model at their levels for this solve, as GAMS's \cmd{x.fx = x.l}:
whole variables, \cmd{hold(P Q)}, or elements, \cmd{hold(P(g1))}, with
\cmd{P(e1,e2)} for two indexes. With \opt{}{from()}, the levels are those of
the simulation.

\item[\optarg{nod}{es}{\#}] and \optarg{}{optcr}{\#} (from version 0.18.0)
set the node limit and the relative gap of branch and bound, which solves a
nonlinear programme with integer or binary variables (\manref{cge_bound}), as
GAMS's \cmd{nodlim} and \cmd{optcr}. The defaults are 10,000 nodes and a gap
of 0, to the tolerance $10^{-9}$.

\item[\optarg{}{starts}{\#}] (from version 0.24.0) solves a nonlinear
programme from \# points, as the multi-start options of KNITRO and BARON
do: the start and $\# - 1$ points of the free variables drawn uniformly
within their bounds, or within $\max(1, |x|)$ of the start where a bound is
missing, keeping the best solution that converges. With integer variables,
every relaxation of branch and bound is started from \# points, so that a
node whose relaxation has a poor local optimum is less likely to be pruned.
The points come from a generator of the package's own and repeat from run
to run; Stata's random numbers are left alone. Multi-start makes a global
optimum of a nonconvex problem more likely, not certain.

\item[\opt{block}{s}] (from version 0.35.0) solves a standalone model
(\manref{cge_model}) block by block: the blocks of the square part of its
system that \manref{cge_structure} finds, in an order in which each depends
only on itself and on those before it, as CONOPT solves a model's
pre-triangular and post-triangular parts apart from its simultaneous core.
Each block is solved by Newton's method on its own equations and variables,
with the exact derivatives of each equation alone, the others held at their
solutions; then the whole system is checked, and polished by Newton's method
if it is not within the tolerance. The solution is that of the whole system:
in the tests, to $10^{-15}$, on a model of five blocks and on a recursive
chain of 600 equations with a simultaneous block of 20, solved in about the
same time either way. A structurally singular system stops with error 459.
It is for standalone models without constraints or inequalities.

\item[\optarg{sol}{ver}{highs}] (from version 0.27.0, an experimental
prototype) solves a linear or mixed-integer linear programme with HiGHS
(Huangfu and Hall 2018), through SciPy's \cmd{linprog} and \cmd{milp} and
Stata's Python integration, in place of the interior point and branch and
bound. The equations and the objective must be linear, which the command
checks at a second point within the bounds (error 459 otherwise). A linear
programme returns the marginals of its equations, in the units of the
interior point's; the statuses follow GAMS's codes, 3 for an unbounded and 4
for an infeasible programme. It works only on Stata 16 and later with Python
and SciPy; Stata 13 to 15 stop with an error.

\item[\optarg{time}{limit}{\#}] (from version 0.16.0) sets a time limit in
seconds, as GAMS's \cmd{reslim}; by default there is none, or the limit set by
\manref{cge_options}. The solvers check it at each iteration, to the second; a
solve that reaches it stops there, with solver status 3.

\item[\optarg{step}{s}{\#}] applies the shocks in \argu{\#} equal steps, each
solution the start of the next. The default, \cmd{steps(0)}, does so only after
a direct solve fails, first in 4 steps and then in 16; \cmd{steps(1)} never
does.

\item[\optarg{min}{imize}{expression}] and \optarg{max}{imize}{expression}
(from version 0.15.0) solve a standalone model as a nonlinear programme: the
free variables minimise or maximise \argu{expression}, an expression of the
user language (\manref{cge_equation}) in the model's variables and
parameters, subject to its equations and to the bounds set by
\manref{cge_bound}. The model must have more free variables than equations,
and its inequalities must be written as equations with slack variables
bounded at zero; a variable with equal bounds should be fixed instead. A
start outside a variable's bounds is moved within them. The tolerance
applies to the scaled conditions for an optimum.

\item[\opt{}{trace}] displays the largest scaled residual, the number of
Jacobians computed and the number of Broyden updates at each iteration.

\item[\opt{diag}{nose}] displays the diagnostics described below after a solve
that converges. They are always displayed after a solve that fails.

\item[\opt{}{full}] solves the full system, with every equation, instead of the
reduced system, from which explicit definitions are left out. Both give the
same solution; \opt{}{full} is for checks, and is much slower for large models
with the dense LU.

\item[\opt{sp}{arse}, \opt{den}{se}] (from version 0.10.1) choose the linear
solver of Newton's steps for this solve, the sparse LU
(section~\ref{sec:sparse}) or the dense one, in place of the setting of
\manref{cge_options}. Both give the same solution.
\end{optdesc}

\cgesec{Remarks and examples}

A simulation is an equilibrium of the model under the closure set by
\manref{cge_closure} and the shocks set by \manref{cge_shock}. \cmd{cge_solve}
finds it by Newton's method, starting from the base, and stores every variable
of the solution with the closure, the shocks and the solution's residual.
Chapter~\ref{chap:solution} describes the method; the remarks below show how
it is used.

\cgesubsec{Example 1: a tariff removal}

The synthetic SAM is mapped and calibrated with uniform elasticities:
\stlog{solve_setup}
Removing every import tariff and solving takes seven iterations:
\stlog{solve_notariff}
The largest scaled residual, over every equation of the model, is below the
default tolerance of \texttt{1e-10}. \cmd{WALRAS} is the slack of the
savings--investment balance, the equation that Walras's law makes redundant
(section~\ref{sec:system}); it is zero, to the precision of the solution, at
every solution. The simulation can now be reported:
\begin{stcode}
. cge_results, sim(notariff)
. cge_report, sims(notariff) tables(macro welfare)
\end{stcode}

\cgesubsec{Example 2: convergence}

With \opt{}{trace}, the residual is displayed at each iteration. One Jacobian
was computed, exactly, and every later step used its LU factors with
Broyden's updates, because each step at least halved the residual:
\stlog{solve_trace}
Newton's method converges quadratically near a solution when the Jacobian is
recomputed at each step. Broyden's updates correct the Jacobian along each
step at the cost of a few inner products, and converge superlinearly; each
step costs one evaluation of the equations and two triangular solves, against
a Jacobian and a factorisation for a Newton step, so for large models they
are much faster. A new Jacobian is computed whenever a step fails to halve
the residual (section~\ref{sec:newton}).

The results are returned in \cmd{r()}:
\stlog{solve_return}
\res{r(N_solved)} is the number of equations in the system that Newton's
method solved, 69 of the model's 93: the others are explicit definitions,
computed from the variables of the reduced system (section~\ref{sec:reduced}).

\cgesubsec{Example 3: a solution that makes no economic sense}

A tenfold rise in government purchases is far outside the range of the base.
With savings-driven investment the equations can still be solved:
\stlog{solve_collapse}
The solution satisfies every equation, but the economy has absorbed the
purchases by negative investment, and the agricultural and manufacturing
activities have all but stopped. \cmd{cge_solve} notes such values, negative
or below one millionth of their base, and stores their number in
\res{r(N_collapsed)}; a solution with collapsed values should be checked
before it is reported.

\cgesubsec{Example 4: a shock the model cannot absorb}

With investment fixed, the same rise cannot be supplied:
\stlog{solve_fail}
The direct solve failed, and so did continuation in 4 and in 16 steps.
\cmd{cge_solve} then stops with error 430 and displays the diagnostics: the
equations with the largest scaled residuals at the last point, by block and
element, and the smallest pivot of the Jacobian relative to the largest, with
the variable it belongs to. A small pivot means that the Jacobian is close to
singular at the last point. Such a failure is usually a property of the shock, the closure or the SAM, not
of the solver: a shock that the economy cannot absorb, a closure with an
equation that no free variable can satisfy, or a SAM with a very small flow
that the shock must expand many times. Equations with no free variable, and
free variables in no equation, are listed when there are any.

Before it solves, \cmd{cge_solve} checks that the system is square, with as
many free variables as equations. If it is not, for example after a variable
added by \cmd{cge_variable, value()} without the equation that determines it,
the command stops with error 459; \manref{cge_variable} shows an example.

\cgesubsec{Constraints}

With constraints set by \manref{cge_constraint}, \cmd{cge_solve} first solves
the complementarity conditions in a smoothed form, which finds whether each
constraint binds, and then solves the model again with the exact conditions
(section~\ref{sec:cons-solution}). The regime of each constraint is displayed
and stored with the simulation, in \res{r(constraints)}.

\cgesubsec{Model and solver statuses}

From version 0.16.0, every solve reports GAMS's model status and solver
status, in \res{r(modelstat)} and \res{r(solvestat)}, and stores them with
the simulation, for \manref{cge_results}. A solve that fails displays them:
\begin{center}
\begin{tabular}{@{}l l p{0.52\textwidth}@{}}
\hline
Model status & Meaning & When\\
\hline
1 & optimal & a solution of a system with constraints\\
2 & locally optimal & a nonlinear programme or an optimisation that converged\\
5 & locally infeasible & the solver stopped without a solution\\
6 & intermediate infeasible & an iteration or time limit stopped the solve\\
7 & feasible solution & a nonlinear programme stopped short at a feasible point\\
8 & integer solution & branch and bound stopped by a limit with an integer solution\\
10 & integer infeasible & branch and bound ended without an integer solution\\
13 & error, no solution & an equation cannot be evaluated at the start\\
14 & no solution returned & branch and bound stopped by a limit without one\\
16 & solved & a solution of a system of equations\\
\hline
Solver status & Meaning & \\
\hline
1 & normal completion & \\
2 & iteration interrupt & the iteration limit was reached\\
3 & resource interrupt & the time limit was reached\\
4 & terminated by solver & the solver stopped without a solution\\
5 & evaluation error limit & an equation cannot be evaluated\\
\hline
\end{tabular}
\end{center}
\manref{cge_options}, Example 3, shows them after iteration and time limits.

\cgesubsec{Example 5: a nonlinear programme}

From version 0.15.0, a standalone model with more free variables than
equations can be solved as a nonlinear programme. Problem 71 of
\citet{hock1981}, which is IPOPT's example, minimises $x_1 x_4 (x_1 + x_2 +
x_3) + x_3$ subject to $x_1 x_2 x_3 x_4 \ge 25$, written with a slack $S \ge
0$, $\sum_i x_i^2 = 40$ and $1 \le x_i \le 5$:
\stlog{solve_nlp}
The optimum, 17.0140173, is the published one. The marginals are the
changes in the objective per unit increase in the right side of each
equation, the marginals of GAMS's solution listing: tightening the product
from 25 raises the minimum by 0.552 per unit, and enlarging the sphere lowers
it by 0.161.

Balancing a matrix to new row and column totals by minimising its cross
entropy against a prior is another nonlinear programme. One column equation
is left out, as the row and column totals have the same sum:
\stlog{solve_entropy}
The solution is the RAS balance of the prior, $X_{rc} = \rho_r A_{rc}
\sigma_c$, which minimises the cross entropy subject to the totals
\citep{mcdougall1999}, to $10^{-11}$, as the test suite checks.
\manref{cge_balance} balances a SAM by these methods directly.

\cgesubsec{Feasibility restoration}

From version 0.36.0, the interior point has a feasibility restoration phase,
after IPOPT's \citep{wachter2006}. Where it cannot find a step, its line
search fails, or the step it takes is too short to move, at a point that
does not satisfy the equations, it looks for a point of smaller
infeasibility: it minimises the squared norm of the scaled residuals by
projected Levenberg--Marquardt steps within the bounds, near the point at
first, until the norm falls by a tenth, and goes on from there with its
slacks and multipliers started again. Where the steps end at a stationary
point of the infeasibility, the programme is locally infeasible (model
status 5), and the method stops there rather than at its iteration limit.
The example of \citet{wachter2000}, on which interior-point methods of this
kind stall at an infeasible point, solves with two restorations:
\stlog{solve_restore}
The residuals are scaled by their base values, so the form of an equation
matters to restoration as to Newton's method. Written as $x_1^2 - x_2 - 1 =
0$ and $x_1 - x_3 - 1/2 = 0$, the scaled residuals have a local minimum near
$x_1 = -1$, toward which the iterates go from the same start: restoration
stops the method there after 10 iterations, locally infeasible, where it
had used all its iterations before. \manref{cge_iis} finds the bounds and
constraints that leave a scenario without a solution.

\cgesubsec{Example 6: integer variables}

From version 0.18.0, integer and binary variables (\manref{cge_bound}) make a
nonlinear programme a mixed-integer one, solved by branch and bound. A
knapsack of eight items, with the capacity written as an equation with a
slack:
\stlog{solve_mip}
The optimum is that of every one of the 256 choices; the nodes are the
relaxations solved.

\cgesubsec{The reduced and the full systems}

Explicit definitions, such as intermediate and household demand, are left out
of the Newton system and computed from the other variables. For a SAM with 250
activities and 20 households, 4,254 of 40,190 equations remain. With
\opt{}{full}, every equation is solved by Newton's method; the solutions agree
to the tolerance, and the test suite checks that they do.

\cgesec{Stored results}

\cmd{cge_solve} stores the following in \cmd{r()}:
\begin{storedres}
\resgroup{Scalars}
\res{r(converged)} & 1 if converged, 0 otherwise\\
\res{r(iterations)} & number of iterations\\
\res{r(max_residual)} & largest scaled residual, over every equation\\
\res{r(N_solved)} & number of equations in the system solved\\
\res{r(jacobians)} & number of Jacobians computed\\
\res{r(broyden)} & number of steps with Broyden's updates\\
\res{r(N_collapsed)} & number of values of a solution that are negative or below one millionth of their base value\\
\res{r(walras)} & value of the slack variable \cmd{WALRAS}\\
\res{r(steps)} & steps of continuation (0 for a direct solve)\\
\res{r(N_constraints)} & number of constraints, if any\\
\res{r(N_binding)} & number of constraints that bind\\
\res{r(modelstat)} & GAMS's model status of the solution (from version 0.16.0)\\
\res{r(solvestat)} & GAMS's solver status\\
\res{r(tolerance)} & the tolerance applied\\
\res{r(iterate)} & the iteration limit applied\\
\res{r(objective)} & the objective at the optimum, for a nonlinear programme (from version 0.15.0)\\
\res{r(N_dof)} & free variables less equations, for a nonlinear programme\\
\resgroup{Macros}
\res{r(name)} & name of the simulation\\
\res{r(method)} & solution method: \cmd{newton}, \cmd{broydenpowell}, with constraints \cmd{semismooth}, \cmd{nlp} or \cmd{minlp}\\
\res{r(objective_expr)} & the objective of a nonlinear programme\\
\res{r(model)} & the named model solved, if any (from version 0.17.0)\\
\res{r(from)} & the simulation the solve started from, if any\\
\res{r(hold)} & the variables held, if any\\
\res{r(N_integer)} & integer and binary elements, for a nonlinear programme (from version 0.18.0)\\
\res{r(nodes)} & nodes of branch and bound\\
\res{r(bestbound)} & the best bound of branch and bound\\
\res{r(gap)} & its relative gap from the objective\\
\res{r(starts)} & starts of the interior point (from version 0.24.0)\\
\res{r(best_start)} & the start of the best solution, with \opt{}{starts} and no integers\\
\res{r(N_starts_converged)} & the starts that converged\\
\res{r(solver)} & \cmd{ipm}, \cmd{bb} or \cmd{highs} (from version 0.27.0)\\
\res{r(hessian)} & the interior point's Hessians: \cmd{exact} or \cmd{differences} (from version 0.28.0)\\
\res{r(N_hessian_exact)} & its iterations with exact Hessians\\
\res{r(N_hessian_differences)} & its iterations with differences\\
\res{r(N_restorations)} & its feasibility restorations (from version 0.36.0)\\
\res{r(linsolver)} & the linear solver: \cmd{dense}, \cmd{sparse} or \cmd{auto} (from version 0.10.1)\\
\res{r(closure)} & closure\\
\res{r(constraints)} & the constraints and their regimes, if any\\
\resgroup{Matrices}
\res{r(marginals)} & the marginals of the equations of a nonlinear programme\\
\resgroup{Diagnostics (after a failure, or with \opt{diag}{nose}; \manref{cge_listing} lists every equation element)}
\res{r(diag_block)} & block of the largest residual\\
\res{r(diag_element)} & its element\\
\res{r(diag_residual)} & that residual\\
\res{r(diag_N_empty_rows)} & number of equations with no free variable\\
\res{r(diag_N_empty_cols)} & number of free variables in no equation\\
\res{r(diag_pivot_ratio)} & smallest pivot of the Jacobian relative to the largest\\
\res{r(diag_singular)} & 1 if the Jacobian is singular or nearly so\\
\end{storedres}

\cgesec{Methods and formulas}

Let $x$ be the free variables of the system and $r(x) = 0$ its equations. The
solver works on scaled variables $z_j = x_j / s_j$ and scaled residuals
$F_i(z) = r_i(x) / e_i$, where $s_j$ and $e_i$ are the base magnitudes of the
variable and of the equation (one where the base value is zero). A Newton step
solves
\begin{equation*}
J(z)\,\Delta z = -F(z),
\end{equation*}
and the new point is $z + \lambda\,\Delta z$, with $\lambda$ the largest of
$1, \tfrac12, \tfrac14, \dots$ such that
\begin{equation*}
\lVert F(z + \lambda\,\Delta z)\rVert \le (1 - 10^{-4}\lambda)\,\lVert F(z)\rVert .
\end{equation*}
After each Newton step, later steps use the LU factors of the last
Jacobian with Broyden's updates \citep{broyden1965, kelley1995}: after a step
$s$, with $y$ the change in $F$, the inverse $H$ of the Jacobian becomes
\begin{equation*}
H_+ = (I + a\, s')\, H, \qquad a = \frac{s - H y}{s' H y},
\end{equation*}
and the next direction is $-H_+ F$, with a line search of at most two
halvings. A new Jacobian is computed when that fails, after a step that
reduces $\lVert F\rVert$ by less than a tenth, and after 30 updates. The solve has converged when $\max_i |F_i| <
\text{\texttt{tolerance()}}$, over every equation of the model, including
those left out of the reduced system; up to three more steps then bring the
residual toward a thousandth of the tolerance.

The Jacobian is exact: the derivatives of every equation are written out,
user equations are differentiated in forward mode, or from version 0.30.0
in reverse mode when an equation has many variables, and in the reduced system
the variables of the left-out definitions $D_k(d_k, y) = 0$ are eliminated
by the chain rule, with $\partial d_k / \partial y_j = -(\partial D_k /
\partial y_j) / (\partial D_k / \partial d_k)$. If Newton fails, it is tried
again with central differences over groups of columns that share no row,
\begin{equation*}
J_{ij} \approx \frac{F_i(z + h \odot e_G) - F_i(z - h \odot e_G)}{2 h_j}, \qquad j \in G,
\qquad h_j = 10^{-6} \max(1, |z_j|),
\end{equation*}
where $e_G$ has ones in the columns of $G$ and $\odot$ is the elementwise
product, then Powell's hybrid method (\cmd{solvenl()}, technique
\cmd{broydenpowell}), followed by Newton steps to polish the result.

Continuation in $n$ steps solves the model for $k = 1, \dots, n$ with
parameters and fixed values
\begin{equation*}
\theta_k = \theta^0 + \frac{k}{n}\,(\theta^1 - \theta^0),
\end{equation*}
between their base values $\theta^0$ and their shocked values $\theta^1$, each
step starting from the solution of the step before.

A nonlinear programme is solved by the primal-dual interior-point method of
\citet{wachter2006} described in section~\ref{sec:nlp}.

Chapter~\ref{chap:solution} gives the details: the reduced system, the
exact Jacobian, Broyden's updates, the grouping of columns for the differences
and the treatment of constraints.

\cgesec{Also see}

\begin{alsosee}
\seeentry{cge_closure}{Choose the model closure}
\seeentry{cge_shock}{Change a parameter or a fixed variable}
\seeentry{cge_constraint}{Limits that bind in some scenarios and not in others}
\seeentry{cge_results}{Report simulation results}
\seeentry{cge_report}{Standard report tables for one or more simulations}
\seeentry{cge_check}{Run model validation checks}
\seechapter{chap:solution}{The solution method}
\end{alsosee}
